\section{引言}
\label{sec:introduction}

不同任务与领域的微调产生了大量专用模型。模型合并通过组合这些模型的参数构建共享模型，降低联合重训的成本，并减少对原始训练数据的依赖~\citep{wortsman2022model,ilharco2023editing}。已有方法通过处理任务增量与参数冲突、调整谱子空间，或利用激活和曲率统计改善合并~\citep{yadav2023ties,marczak2025no,jin2023regmean,tam2024mats}。然而，合并模型与各任务专家之间仍存在性能差距。

对于基于统计匹配的合并，构造代理目标与求解该目标是两个环节。MaTS 使用共轭梯度求解由任务统计确定的线性系统~\citep{tam2024mats}。有限校准样本可能使系统低秩或病态，直接求逆会放大估计误差，过强的正则化又会削弱任务信息。即使初始系统被充分求解，它所描述的仍是构造时的模型状态：更新前层会改变后层输入，预测改善也会改变尚待修正的专家分歧。因此，后续更新所需的统计可能已经变化，继续求解初始代理未必能反映当前需求。

合并过程中的表示变化已有研究关注。RegMean++ 将合并模型的层内与跨层依赖纳入回归目标~\citep{nguyen2026regmeanpp}；IterIS 在 LoRA 合并中交替推断特征与求解，逐步修正优化目标~\citep{chen2025iteris}。本文关注多个固定任务专家与当前共享模型之间的剩余功能分歧，并在每次接受更新后重新测量相应的修正需求。

我们提出 \textbf{\method{}（CM）}，将合并组织为测量、求解、选步与重测的循环。每轮在少量无标签校准数据上，以固定专家为教师，收集当前模型的层输入统计与教师学生 KL 修正统计；再以专家参数为锚点，用有限步迭代求解各层候选权重。候选模型由留出集上的教师学生 KL 选择更新幅度，接受后重新收集统计。数值稳定化、更新幅度选择与外层重测分别处理求解可靠性、整网更新和统计适配问题。

机制实验支持这一设计：修正度量在早期更新后明显变化，随后逐渐趋稳。在 Average 起点的配对实验中，沿用初始统计仍有收益，但低于每轮重测；将单轮求解增加到 400 步，也未达到四轮重测的结果。CLIP、T5 与 LLM 实验进一步表明，CM 可以改善多种合并起点，并在部分设定中缩小起点间的差距。收益随任务和初始化而异，完整结果保留了负向案例。

本文的主要贡献如下：
\begin{itemize}
  \setlength{\itemsep}{2pt}
  \item 区分固定代理的数值求解与更新后的统计重测，通过配对实验检验两者对合并效果的作用。
  \item 提出 CM，以当前模型的剩余专家分歧构造逐层修正，并结合留出集选步与接受后的统计重测。
  \item 在视觉、自然语言处理和大语言模型上验证 CM，系统报告主配置 CM + Iso-C，以及从 WeightAvg、RegMean 等方法开始的起点适配结果与边界。
\end{itemize}
